% Демонстрация 1: «Тройное произведение = объём» (Лекция 8)
\section*{Демонстрация 1. «Тройное произведение = объём»\\ \normalsize\textmd{Файл \texttt{demos/01-triple-product.html} \quad(к Лекции~8)}}

\subsection*{Какие факты лекции иллюстрирует}

Демонстрация представляет собой интерактивную 3D-сцену (Three.js). Пользователь
ползунками задаёт координаты трёх векторов $\va=(a_x,a_y,a_z)$,
$\vb=(b_x,b_y,b_z)$, $\vc=(c_x,c_y,c_z)$ (каждая компонента меняется в пределах
от $-3$ до $3$ с шагом $0{,}1$). На сцене в реальном времени отображаются:
три цветные стрелки-вектора $\va$, $\vb$, $\vc$ из начала координат;
полупрозрачный параллелепипед, построенный на этих векторах; вектор
$[\vb,\vc]$ (включается галочкой); оси координат и подписи. В боковой панели
синхронно пересчитываются три числа: смешанное произведение
$(\va,\vb,\vc)=\det$, объём $V=|(\va,\vb,\vc)|$ и ориентация тройки (правая /
левая), а при вырождении загорается плашка «Векторы компланарны! Объём $=0$».
Сцену можно вращать, приближать и перемещать мышью.

Тем самым демонстрация наглядно показывает следующие конкретные факты
Лекции~8.

\begin{itemize}
  \item \textbf{Смешанное произведение как ориентированный объём}
    (теорема о геометрическом смысле, см. рисунок~8.4): число $(\va,\vb,\vc)$,
    показанное в панели, в точности совпадает с объёмом параллелепипеда на
    сцене, взятым со знаком. При правой тройке $(\va,\vb,\vc)=+V$, при левой
    $(\va,\vb,\vc)=-V$. Студент видит, что меняется именно \emph{знак}, а не
    форма тела.
  \item \textbf{Определение смешанного произведения}
    $(\va,\vb,\vc)=([\va,\vb],\vc)$ и его координатная запись через
    определитель $3\times3$ (формула~(8.7)): подпись панели «$(\va,\vb,\vc)=
    \det$» прямо связывает геометрическую картинку с вычислением
    \[
      (\va,\vb,\vc)=
      \begin{vmatrix}
        a_x & a_y & a_z \\
        b_x & b_y & b_z \\
        c_x & c_y & c_z
      \end{vmatrix}.
    \]
  \item \textbf{Критерий компланарности} (свойство~3):
    $(\va,\vb,\vc)=0 \Leftrightarrow \va,\vb,\vc$ компланарны. Когда тройка
    приближается к компланарной, параллелепипед \emph{схлопывается} в плоскую
    фигуру, число $V$ стремится к нулю, а грани становятся почти прозрачными
    (прозрачность граней в демонстрации масштабируется по объёму). Это —
    зрительный образ вырождения определителя.
  \item \textbf{Ориентация упорядоченной тройки} (определение правой и левой
    тройки): индикатор «Правая\,/\,Левая» переключается ровно в тот момент,
    когда знак определителя меняется. Поменяв местами два вектора
    (свойство о смене знака при транспозиции, $(\va,\vb,\vc)=-(\vb,\va,\vc)$),
    студент видит смену ориентации при сохранении объёма.
  \item \textbf{Связь со скалярным и векторным произведениями}: галочка
    «Показать $[\vb,\vc]$» выводит вектор векторного произведения. Можно
    наблюдать, что $(\va,\vb,\vc)=([\vb,\vc],\va)$ (циклическая симметрия,
    свойство~1) есть просто проекция $\va$ на $[\vb,\vc]$, умноженная на длину
    $[\vb,\vc]$, то есть «площадь основания $\times$ высота».
  \item \textbf{Циклическая симметрия} $(\va,\vb,\vc)=(\vb,\vc,\va)=
    (\vc,\va,\vb)$ (свойство~1): циклически переставляя значения ползунков,
    студент убеждается, что и число, и объём, и ориентация не меняются —
    в полном соответствии с примерами лекции на циклическую симметрию.
\end{itemize}

\begin{remarkIT}{Что из этого «живёт» в реальных приложениях}{demo01_apps}
  Демонстрация даёт зрительную опору сразу нескольким прикладным сюжетам
  лекции:
  \begin{itemize}
    \item \textbf{Компьютерная графика:} вектор $[\vb,\vc]$ на сцене — это
      нормаль к грани (как при импорте 3D-меша), а знак $(\va,\vb,\vc)$
      определяет, с какой стороны грань «лицевая» (back-face culling).
    \item \textbf{Проверка точки внутри многогранника:} наблюдаемая смена
      знака определителя при переходе вектора через плоскость основания —
      это и есть основа теста принадлежности точки выпуклому телу по знакам
      смешанных произведений.
    \item \textbf{3D-реконструкция (SfM, SLAM):} схлопывание параллелепипеда
      при $(\va,\vb,\vc)=0$ — наглядная модель вырожденной (компланарной)
      конфигурации точек, непригодной для триангуляции глубины.
    \item \textbf{Робототехника:} $V=|(\va,\vb,\vc)|$ — оценка объёма рабочей
      зоны манипулятора, а компланарность сигнализирует о вырожденной
      (особой) конфигурации суставов.
  \end{itemize}
\end{remarkIT}

\subsection*{Примерный сценарий использования на занятии}

Демонстрацию рекомендуется включать \emph{после} введения определения
смешанного произведения~(8.6) и \emph{во время} разбора теоремы о его
геометрическом смысле. Ниже — примерный пошаговый сценарий (около
12--15~минут).

\begin{enumerate}
  \item \textbf{Запуск и калибровка взгляда (1--2~мин).}
    Откройте демонстрацию со стартовыми значениями
    $\va=(2,0,0)$, $\vb=(0,2,0)$, $\vc=(0,0,2)$ — это ортогональная правая
    тройка, привычный «угол комнаты». Покажите, что в панели $\det=8$,
    $V=8$, ориентация «Правая». Спросите аудиторию: \emph{«Почему именно
    $8$?»} (ответ: $2\cdot2\cdot2$ — объём куба со стороной~2). Поверните
    сцену мышью, чтобы все увидели параллелепипед с трёх сторон.

  \item \textbf{Связь числа и определителя (1--2~мин).}
    Обратите внимание на формулу в нижней панели и запишите на доске тот же
    определитель~(8.7) с текущими координатами. Подчеркните: \emph{число в
    панели — это не «ещё одна» величина, а значение того самого определителя,
    что мы вычисляем на бумаге}.

  \item \textbf{Деформация без потери объёма (2~мин).}
    Сдвиньте ползунок $c_x$ или $c_y$ так, чтобы вектор $\vc$ наклонился, но
    высота параллелепипеда почти не изменилась. Покажите: тело «перекосилось»,
    но $V$ держится почти на месте. Задайте вопрос: \emph{«От чего зависит
    объём — от наклона рёбер или от высоты над основанием?»} Свяжите с
    выкладкой $(\va,\vb,\vc)=S_{\text{осн}}\cdot|\vc|\cos\varphi=S_{\text{осн}}
    \cdot h$ из доказательства (рисунок~8.4).

  \item \textbf{Момент удивления: схлопывание (2--3~мин).}
    Медленно ведите ползунок $c_z$ к нулю (оставив $c_x,c_y$ маленькими),
    пока вектор $\vc$ не ляжет в плоскость $\va,\vb$. Параллелепипед на глазах
    \emph{сплющивается в плоскую фигуру}, число $V$ скатывается к нулю и
    загорается красная плашка «Векторы компланарны! Объём $=0$». Это и есть
    кульминация: студент \emph{видит}, как $(\va,\vb,\vc)=0$ означает
    компланарность (свойство~3). Спросите: \emph{«Что стало с высотой? А с
    определителем?»}

  \item \textbf{Возврат и смена ориентации (2~мин).}
    Верните $\vc$ наверх, затем проведите его \emph{вниз}, на отрицательные
    $z$. Объём восстановится, но индикатор переключится на «Левая», а
    определитель станет отрицательным. Зафиксируйте формулировку теоремы:
    $(\va,\vb,\vc)=\pm V$, знак кодирует ориентацию. Хороший момент спросить:
    \emph{«Поменялось ли тело? А его объём?»} (нет; поменялся только знак).

  \item \textbf{Транспозиция и циклическая симметрия (2~мин).}
    Поменяйте местами значения ползунков $\va$ и $\vb$ (или просто переставьте
    тройку циклически). Покажите на числах: при транспозиции знак меняется
    $(\va,\vb,\vc)=-(\vb,\va,\vc)$, а при циклическом сдвиге всё сохраняется
    $(\va,\vb,\vc)=(\vb,\vc,\va)$ (свойства~1 и~2). Это «живая» версия
    разобранного в лекции примера на свойства смешанного произведения.

  \item \textbf{Связь с векторным произведением (1--2~мин).}
    Включите галочку «Показать $[\vb,\vc]$». Появится фиолетовая стрелка,
    перпендикулярная грани на $\vb,\vc$. Объясните: $(\va,\vb,\vc)=
    ([\vb,\vc],\va)$ — проекция $\va$ на нормаль, умноженная на площадь
    основания $|[\vb,\vc]|$. Когда $\va$ почти лежит в плоскости основания
    ($\va\perp[\vb,\vc]$), проекция мала — и объём мал. Свяжите с критерием
    компланарности и приложением «нормаль к полигону» из \texttt{remarkIT}.

  \item \textbf{Переход к практике (1~мин).}
    Завершите так: \emph{«Теперь, не глядя на картинку, по одним координатам
    предскажите знак и объём»} — и переходите к разбору примеров лекции:
    объём параллелепипеда, объём пирамиды (множитель
    $\tfrac16$) и компланарность четырёх точек. Демонстрация
    остаётся открытой: координаты из условий примеров можно ввести ползунками
    и сверить знак и объём с вычислениями на доске.
\end{enumerate}

\noindent\textbf{Методический совет.} Главный дидактический эффект даёт
именно \emph{переход через нуль}: пусть студенты сами назовут момент, когда
объём обращается в нуль, и только потом запишите критерий
$(\va,\vb,\vc)=0\Leftrightarrow$ компланарность. Демонстрация превращает
абстрактное «определитель равен нулю» в наглядное «коробка сплющилась».
